\documentclass[a4paper,12pt]{article}
\usepackage[left=2.5cm,right=2.5cm,top=0cm]{geometry}
\usepackage{tikz}

\usetikzlibrary{shapes.geometric}
\usepackage{hyperref}
\hypersetup{
    colorlinks=false,
    urlbordercolor=white
    }
\usepackage[none]{hyphenat}
\title{Monthly Maths Circle India Challenge}
% \author{Disha Kuzhively}
% \email{disha.jk@icts.res.in}
\date{Due Date: 29 February 2024}
\begin{document}
\maketitle
\thispagestyle{empty}
\begin{enumerate}
    \item[Problem 1.] Anita and Bindu collaborated to create a festive cake arrangement. Anita contributed four cakes, each an isosceles triangle, while Bindu baked a single large cake shaped like a trapezium. When they assembled the cakes at the party, a delightful surprise awaited them - the arrangement formed a regular five-pointed star, as in the picture below.
    \begin{figure}[h]
        \centering
        \begin{tikzpicture}
    \def\radius{1cm}
    \def\factor{1/cos(72)}
    \def\infactor{1/cos(36)}
\node [name=s, shape=star, star points=5, star point height=\radius*(\factor-\infactor), minimum size=2*\radius*\factor, draw,very thick] at (0,0) {};
\filldraw[fill=gray!50,very thick, draw=black] (s.inner point 2) -- (s.inner point 4) -- (s.outer point 5) -- (s.outer point 2) -- cycle;
\end{tikzpicture}
                
    \end{figure}

    A regular pentagonal star is created by connecting alternate vertices of a regular pentagon.

    Now comes the time to split the cost. Bindu claims that the total area of Anita's four triangular cakes equals the single large trapezium cake. Anita argues that the combined area of her cakes might be larger. Can you help them split the cost?

    % A pentagram is a regular five-pointed star formed by connecting alternate vertices of a pentagon to create the points of the star. Show that the area of the shaded region of the regular pentagonal star in the picture below is exactly half of the total area of the star.
    
    \item[Problem 2.] Construct a square on a square grid with an area five times that of a single square on the grid using only a straight edge and a pen or pencil.
    \begin{figure}[h]
        \centering
        \begin{tikzpicture}
    \draw[gray] (-1.5,-1.5) grid (2.5,2.5);
    \draw[black,thick] (0,0) grid (1,1);
    \draw[->] (0.5,0.5) -- (3,1.5) node[above right] {single square on the grid};
    \end{tikzpicture}
    \end{figure}
    % A group of friends decided to buy plots together
\end{enumerate}
\newpage
\thispagestyle{empty}
\centering
\begin{minipage}{0.8\textwidth}
    \vspace{3cm}
    Monthly Maths Circle India Challenge is a maths competition open to everyone who likes to solve mathematical problems. Every second Friday of the month, we will upload a new problem on our website. You will have three weeks to solve it and submit your answer before the deadline specified on the challenge sheet. Anyone can join, regardless of their age. We seek creative solutions that show a good understanding of various mathematical concepts. Please avoid using computer programs or excessively lengthy calculations. We will evaluate your submission based on its completeness, accuracy, and presentation quality. Remember that this contest aims to stimulate mathematical creativity; our graders' decisions are final.
\end{minipage}
\centering
\begin{minipage}{0.8\textwidth}
    \vspace{3cm}
Check \url{https://www.icts.res.in/outreach/mci} for more information
\end{minipage}
\begin{minipage}{0.8\textwidth}
    \vspace{3cm}
    \centering
    \includegraphics[scale=0.15]{mci-monthly}

    Scan to submit your entry.
\end{minipage}
\end{document}